\documentclass[10pt]{article}
\usepackage[T1]{fontenc}
\usepackage[utf8]{inputenc}
\usepackage[french]{babel}
\usepackage{amsmath,amssymb,mathrsfs,mathpazo}
\usepackage{geometry}
\geometry{a4paper,margin=2cm}
\usepackage{graphicx,array,tabularx,tabularray,colortbl,multirow,multicol,booktabs,enumitem,pifont,ulem,fancybox,diagbox,cancel}
\usepackage{pstricks,pstricks-add,pst-plot,pst-tree,pst-node,pst-text,pst-func,pst-math,pst-bspline,pst-3dplot,pst-eucl,pst-solides3d}
\let\pscalculate\undefined
\usepackage{pst-bezier}
\usepackage{tikz,pgfplots,tkz-tab}
\usetikzlibrary{arrows,arrows.meta,patterns,calc,positioning,trees,decorations.markings,shapes.geometric,shapes.misc,shapes.arrows,intersections,decorations.pathmorphing,decorations.pathreplacing,backgrounds,shadows,babel}
\pgfplotsset{compat=1.18}
\usepgfplotslibrary{fillbetween}
\usepackage[np]{numprint}
\usepackage{listings,makecell,pas-tableur,twoopt}
\let\barre\undefined
\usepackage[pstricks]{tablvar}
\definecolor{darkgreen}{RGB}{0,100,0}
\providecommand{\red}{\color{red}}
\providecommand{\blue}{\color{blue}}
\providecommand{\green}{\color{green}}
\providecommand{\black}{\color{black}}
\providecommand{\white}{\color{white}}
\providecommand{\gray}{\color{gray}}
\providecommand{\cyan}{\color{cyan}}
\providecommand{\magenta}{\color{magenta}}
\providecommand{\yellow}{\color{yellow}}
\providecommand{\violet}{\color{violet}}
\providecommand{\orange}{\color{orange}}
\providecommand{\euro}{\texteuro}
\providecommand{\EUR}{\texteuro}
\providecommand{\ieme}{\textsuperscript{e}}
\providecommand{\degres}{\ensuremath{^\circ}}
\providecommand{\degree}{\ensuremath{^\circ}}
\providecommand{\vv}[1]{\overrightarrow{#1}}
\providecommand{\Coord}[1]{#1}
\providecommand{\decosix}{}
\providecommand{\starredbullet}{\ensuremath{\bullet}}
\pagestyle{empty}
\setlength{\parindent}{0pt}
\usepackage{adjustbox}\DeclareUnicodeCharacter{2236}{:}\DeclareUnicodeCharacter{2190}{\ensuremath{\leftarrow}}
\providecommand{\R}{\mathbb{R}}
\providecommand{\N}{\mathbb{N}}
\providecommand{\Z}{\mathbb{Z}}
\providecommand{\Q}{\mathbb{Q}}
\providecommand{\D}{\mathbb{D}}
\providecommand{\C}{\mathbb{C}}
\providecommand{\e}{{\mathrm{e}}}
\providecommand{\vect}[1]{\overrightarrow{#1}}
\providecommand{\vectt}[1]{\overrightarrow{\mathrm{#1}}}
\providecommand{\vvt}[1]{\overrightarrow{\mathrm{#1}}}
\providecommand{\np}[2][]{{#2~#1}}
\providecommand{\vv}[1]{{\overrightarrow{#1}}}
\renewcommand{\d}{{\mathrm{d}}}
\providecommand{\degres}{{^\circ}}
\providecommand{\degre}{{^\circ}}
\providecommand{\blue}{\color{blue}}
\providecommand{\red}{\color{red}}
\providecommand{\enskip}{\quad}
\providecommand{\llbracket}{\lbrack\!\lbrack}
\providecommand{\rrbracket}{\rbrack\!\rbrack}
\providecommand{\euro}{{\text{€}}}
\providecommand{\no}{{\text{n°}}}
\providecommand{\interval}[2]{{[#1;#2]}}
\providecommand{\intervFO}[2]{{[#1;#2[}}
\providecommand{\dotfill}{\ldots}

\providecommand{\ds}{\displaystyle}
\providecommand{\textding}[1]{#1}
\providecommand{\Oij}{\(\mathrm{O};\vec{\imath},\vec{\jmath}\)}
\providecommand{\Oijk}{\(\mathrm{O};\vec{\imath},\vec{\jmath},\vec{k}\)}
\providecommand{\Ouv}{\(\mathrm{O};\vec{u},\vec{v}\)}
\providecommand{\intervFF}[2]{\ensuremath{[#1;#2]}}
\renewcommand{\euro}{\ifmmode\text{\texteuro}\else\texteuro\fi}
\DeclareUnicodeCharacter{00B0}{\ifmmode^{\circ}\else\textdegree\fi}
\DeclareUnicodeCharacter{20AC}{\ifmmode\text{\texteuro}\else\texteuro\fi}
% Énoncés : Éducation nationale. Transcription et corrigés : APMEP, Denis Vergès et contributeurs. Découpage et classement : CoopMaths, Sébastien Lozano, Éric Elter et contributeurs. Adaptation de la présentation : Maths Explorer.
% CC BY-SA : https://coopmaths.fr/www/about/
% https://coopmaths.fr/alea/static/bac/2022/tex/bac_2022_09_sujet1_polynesie_3.tex
\begin{document}
\noindent\textit{APMEP / CoopMaths. CC BY-SA. Présentation adaptée par Maths Explorer.}\par\medskip
\bigskip

\textbf{Partie 1}

\medskip

Soit $g$ la fonction définie pour tout nombre réel $x$ de l'intervalle $]0~;~+ \infty[$ par:

\[g(x) = \dfrac{2 \ln x}{ x}.\]

\smallskip

\begin{enumerate}
\item On note $g'$ la dérivée de $g$. Démontrer que pour tout réel $x$ strictement positif : 

\[g'(x) = \dfrac{2 - 2\ln x}{x^2}.\]

\item On dispose de ce tableau de variations de la fonction $g$ sur l'intervalle $]0~;~+ \infty[$ :

\begin{center}
\psset{unit=1cm}
\begin{pspicture}(9,3)

\psframe(9,3)\psline(0,2.5)(9,2.5)
\psline(2.5,0)(2.5,3)
\uput[u](1.25,2.4){$x$}\uput[u](3,2.4){$0$}\uput[u](5.5,2.4){1}\uput[u](7,2.4){e}\uput[u](8.5,2.4){$+ \infty$}
\psline(2.95,0)(2.95,2.5)\psline(3.05,0)(3.05,2.5)
\uput[u](3.5,0){$- \infty$}\rput(5.5,1.25){0}\uput[d](7,2.5){$\frac{2}{\text{e}}$}\uput[u](8.8,0){0}
\rput(1.25,1.5){Variations}
\rput(1.25,1.){de $g$}
\psline(4,0.32)(5.3,1.1)\psline{->}(5.6,1.3)(6.8,2.)\psline{->}(7.2,2)(8.6,0.4)
\end{pspicture}
\end{center}

\smallskip

Justifier les informations suivantes lues dans ce tableau:

	\begin{enumerate}
		\item la valeur $\dfrac{2}{\text{e}}$ ;
		\item les variations de la fonction $g$ sur son ensemble de définition ;
		\item les limites de la fonction $g$ aux bornes de son ensemble de définition.
	\end{enumerate}
\item En déduire le tableau de signes de la fonction $g$ sur l'intervalle $]0~;~+ \infty[$.
\end{enumerate}

\bigskip

\textbf{Partie 2}

\medskip

Soit $f$ la fonction définie sur l'intervalle $]0~;~+ \infty[$ par

\[f(x) = [\ln (x)]^2.\]

 Dans cette partie, chaque étude est effectuée sur l'intervalle $]0~;~+ \infty[$.

\medskip

\begin{enumerate}
\item Démontrer que sur l'intervalle $]0~;~+ \infty[$, la fonction $f$ est une primitive de la fonction $g$.
\item À l'aide de la partie 1, étudier :
	\begin{enumerate}
		\item la convexité de la fonction $f$ ;
		\item les variations de la fonction $f$.
\end{enumerate}
\item 
	\begin{enumerate}
		\item Donner une équation de la tangente à la courbe représentative de la fonction $f$ au point d'abscisse e.
		\item En déduire que, pour tout réel $x$ dans $]0~;~\text{e}]$ :
		
\[[\ln (x)]^2 \geqslant \dfrac{2}{\text{e}}x - 1.\]

	\end{enumerate}
\end{enumerate}
\bigskip

\end{document}
